\documentclass[11pt,a4paper]{article}
\usepackage[utf8]{inputenc}
\usepackage[italian]{babel}
\usepackage[T1]{fontenc}
\usepackage{cmap}
\usepackage[margin=2cm]{geometry}
\usepackage[table]{xcolor}
\usepackage{enumitem}
\usepackage{hyperref}
\usepackage{tcolorbox}
\usepackage{titlesec}
\usepackage{booktabs}
\usepackage{tikz}
\usepackage{float}
\usepackage{amsmath}
\usepackage{amssymb}

\definecolor{primaryColor}{HTML}{0054A6}
\definecolor{secondaryColor}{HTML}{4A5568}
\definecolor{backgroundColor}{HTML}{F7FAFC}
\definecolor{borderColor}{HTML}{E2E8F0}
\definecolor{successColor}{HTML}{38A169}
\definecolor{warningColor}{HTML}{DD6B20}

\hypersetup{
    colorlinks=true,
    linkcolor=primaryColor,
    urlcolor=primaryColor,
}

\titleformat{\section}
  {\color{primaryColor}\normalfont\Large\bfseries}
  {}{0em}{}
  [\color{borderColor}\titlerule]

\titleformat{\subsection}
  {\color{primaryColor}\normalfont\large\bfseries}
  {}{0em}{}

\titleformat{\subsubsection}[runin]
  {\normalfont\bfseries\color{secondaryColor}}
  {}{0em}{}[:]

\begin{document}

\begin{tcolorbox}[colback=primaryColor,colframe=primaryColor,arc=3mm,halign=center,valign=center,height=2.5cm]
    \color{white}\LARGE\bfseries FINAL REFLECTIONS \\[0.2cm]
    \color{white!80}\large UX Design Project: Help People Help People --- TPER S.p.A. \\[0.1cm]
    \color{white!60}\normalsize Phase C --- Design Proposal
\end{tcolorbox}

\vspace{0.5cm}

\section{Come le Personas Hanno Influenzato il Design}

Le quattro personas sviluppate nel Cast Card hanno guidato ogni decisione progettuale della Phase C. Un tema trasversale emerso dai test \`e la \textbf{confusione sui canali di acquisto}: l'ecosistema TPER si articola su molteplici canali (solweb, Roger, Punti Tper, tabaccherie, contactless a bordo), ma il sito informativo \texttt{tper.it} non forniva alcuna guida per orientarsi. Di seguito viene descritto l'impatto specifico di ciascuna persona sul design proposto.

\subsection{Impatto di Maria Borriello (79 anni, competenza digitale minima, SUS 27.5)}

Maria rappresenta il caso limite del progetto. I suoi 0 task su 5 completati e il tempo medio di 652 secondi per task hanno imposto vincoli progettuali non negoziabili. La difficolt\`a principale di Maria non era l'acquisto in s\'e, ma capire \textbf{dove} andare per acquistare:

\begin{itemize}
    \item \textbf{Card canali con icone visive}: Maria non riusciva a distinguere tra ``solweb'', ``Roger'' e ``Punto Tper'' perch\'e nomi tecnici senza significato intuitivo. Il redesign introduce card visuali con icone grandi e colori distintivi: una carta con il simbolo del telefono per Roger, una carta con il simbolo del computer per solweb, una carta con il simbolo del negozio per Punti Tper. Ogni card mostra solo le 2--3 operazioni principali in linguaggio semplicissimo (``Compra biglietto'', ``Ricarica'', ``Abbonamento'').
    \item \textbf{Tour guidato ``Prima Volta''}: Maria alla prima visita vede una finestra introduttiva opzionale: ``Ti aiutiamo a orientarti'' con 4 schermate illustrate che spiegano i canali con metafore quotidiane (``Come una tabaccheria digitale'', ``Come un negozio fisico'').
    \item \textbf{Testo grande e contrasto elevato (R8)}: La difficolt\`a di lettura riportata nei test ha imposto carattere minimo 14pt e contrasto 4.5:1, particolarmente critico per leggere i nomi dei canali.
    \item \textbf{Treno incoraggiante (R9)}: Maria abbandonava quando incontrava un canale non familiare. I nuovi messaggi la guidano: ``Non preoccuparti, ti spiego dove andare''.
    \item \textbf{Progress bar (R3)}: Nel wizard ``Dove Comprare'', Maria vede sempre a che punto del percorso di scelta si trova (``Passo 2 di 3: scegli come pagare'').
\end{itemize}

\subsection{Impatto di Laura Palumbo (70 anni, competenza media, 7 errori)}

Laura ha dimostrato che anche pochi anni di differenza anagrafica e un uso minimo dello smartphone producono un gap prestazionale significativo. Laura possedeva competenze digitali di base ma \`e stata sopraffatta dalla molteplicit\`a di opzioni:

\begin{itemize}
    \item \textbf{Wizard ``Dove vuoi comprare?''}: Laura vedeva 5+ canali e non sapeva quale scegliere. Il wizard riduce la scelta a un percorso guidato di 2--3 domande semplici: ``Cosa ti serve?'' $\rightarrow$ ``Come vuoi pagare?'' $\rightarrow$ ``Ecco il canale giusto per te''. Ogni risposta riduce le opzioni visibili fino a mostrare solo il canale consigliato.
    \item \textbf{Tabella comparativa dei canali}: Laura faceva confronti manuali aprendo 5 schede. La tabella ``Canali a confronto'' mostra in un colpo d'occhio: quali operazioni supporta ogni canale, costi, orari, lingua. Filtrabile per operazione desiderata (R2).
    \item \textbf{Separazione contenuti informativi/transazionali}: Laura si perdeva tra news aziendali e azioni concrete. La homepage \`e stata riorganizzata per priorit\`a d'uso (R1), con la sezione ``Dove Comprare'' sempre visibile e le news in sezione secondaria.
    \item \textbf{Etichette chiare (R2)}: Laura non capiva ``solweb'' e ``MyAtac'' come nomi di canali. Il redesign usa etichette descrittive: ``Abbonamenti online (solweb)'', ``Biglietti digitali (Roger)'' con il nome tecnico tra parentesi come supporto.
\end{itemize}

\subsection{Impatto di Chris Harper (34 anni, competenza alta, 226 secondi Task 1)}

Chris ha dimostrato che le criticit\`a non colpiscono solo utenti fragili. Nonostante l'elevata competenza digitale, Chris ha perso tempo prezioso a causa della frammentazione dei canali:

\begin{itemize}
    \item \textbf{Confronto canali upfront}: Chris ha impiegato 226 secondi per orientarsi tra i vari canali, visitando 3 domini diversi prima di trovare quello giusto. Il redesign mostra una tabella comparativa dei canali gi\`a nella homepage, permettendo a Chris di scegliere il canale corretto al primo tentativo, senza navigazioni esplorative.
    \item \textbf{Transizione dominio segnalata (UR-01)}: Chris passava da \texttt{tper.it} a \texttt{solweb.tper.it} senza preavviso e doveva riorientarsi. Il redesign mostra un avviso esplicito prima di ogni reindirizzamento: ``Stai per andare su \texttt{solweb.tper.it} per completare l'acquisto dell'abbonamento'' con un pulsante ``Continua'' e uno ``Torna indietro''.
    \item \textbf{Cross-link ``Intendevi...?''}: Chris finiva spesso sul canale sbagliato perch\'e i nomi erano ambigui. Il sistema ora riconosce l'intenzione: se l'utente cerca ``abbonamento'' su una pagina informativa di \texttt{tper.it}, compare un link diretto: ``Stai cercando di acquistare un abbonamento? Vai su \texttt{solweb.tper.it}'' (UR-03).
    \item \textbf{Checklist preventiva (R11)}: Chris doveva cercare autonomamente quali documenti servivano per l'abbonamento. Il redesign mostra una checklist visiva ``Documenti necessari'' gi\`a nella guida informativa, prima del reindirizzamento a solweb.
\end{itemize}

\subsection{Impatto di Fatima Hassan (45 anni, lingua araba)}

Fatima ha evidenziato le barriere linguistiche nell'orientamento tra i canali, non solo nell'uso dei servizi:

\begin{itemize}
    \item \textbf{Guida ai canali multilingua (R12)}: La guida ``Dove Comprare'' \`e tradotta integralmente in arabo e inglese, incluse le descrizioni dei canali, i nomi delle operazioni e i messaggi del wizard. Fatima pu\`o scegliere il canale giusto nella sua lingua senza dover interpretare termini italiani.
    \item \textbf{Ricerca con sinonimi in pi\`u lingue}: Il motore di ricerca accetta ``biglietto'', ``ticket'' e ``tadhkira'' (biglietto in arabo) e restituisce gli stessi risultati, con indicazione del canale appropriato per ogni tipo di acquisto.
    \item \textbf{Segnaletica visiva universale}: Le icone dei canali sono state progettate per essere culturalmente neutre e universalmente riconoscibili: telefono per Roger, edificio per Punto Tper, computer per solweb, carta di credito per contactless. Riducono la dipendenza dalla lingua scritta.
    \item \textbf{QR code multilingua}: Il QR code sullo scontrino del Punto Tper atterra su una pagina informativa che rileva la lingua del dispositivo e si presenta automaticamente in arabo, inglese o italiano.
\end{itemize}

\vspace{0.3cm}

\section{Come gli Use Case Hanno Influenzato il Design}

Gli Use Case (UR-01--UR-05) emersi dal Cast Card hanno guidato le scelte progettuali, con particolare attenzione alla frammentazione dei canali e alla confusione che ne deriva:

\begin{itemize}
    \item \textbf{UR-01 (Frammentazione dei canali)}: Ha portato alla progettazione di \texttt{tper.it} come \textbf{hub informativo integrato}. Invece di isolare ogni canale, il redesign presenta una vista unificata dell'ecosistema TPER, con una scheda per ogni canale, le operazioni supportate e link diretti. L'avviso di transizione esplicito prepara l'utente al cambio di contesto.
    \item \textbf{UR-02 (Confusione post-acquisto)}: Ha generato la sezione ``Dopo l'acquisto'' per ogni canale. Su \texttt{tper.it}, ogni pagina di prodotto include una guida post-acquisto: ``Hai appena comprato su solweb? Ecco come attivare l'abbonamento'' oppure ``Hai ricaricato in tabaccheria? Verifica il credito residuo su Roger'' (R6).
    \item \textbf{UR-03 (Errori dovuti al canale sbagliato)}: Ha portato all'implementazione dei \textbf{cross-link intelligenti}. Quando un utente visita una pagina informativa su \texttt{tper.it} che descrive un prodotto acquistabile altrove, il sistema mostra un suggerimento contestuale: ``Intendevi acquistare? Vai su [canale corretto]''. Questo riduce gli errori di navigazione e gli abbandoni (R10).
    \item \textbf{UR-04 (Ponte fisico/digitale)}: Ha ispirato il sistema di \textbf{QR code} che collega l'esperienza fisica (Punto Tper) a quella informativa (\texttt{tper.it}). Lo scontrino stampato allo sportello include un QR code che porta alla pagina informativa dedicata, con riepilogo dell'operazione, documenti allegati e possibilit\`a di proseguire digitalmente. Il ponte \`e bidirezionale: dalla pagina informativa si pu\`o prenotare un appuntamento al Punto Tper pi\`u vicino.
    \item \textbf{UR-05 (Guida all'abbonamento)}: Ha generato la guida passo-passo ``Come abbonarsi'' in 5 fasi, interamente ospitata su \texttt{tper.it} come contenuto informativo. Ogni fase spiega cosa serve, quali documenti preparare e, al termine, reindirizza l'utente al canale corretto (solweb per acquisto online, Punto Tper per acquisto fisico). Include un calcolatore ``Hai diritto a uno sconto?'' che valuta l'eleggibilit\`a dell'utente in meno di un minuto.
\end{itemize}

\vspace{0.3cm}

\section{Riflessioni sul Processo Progettuale}

\subsection{Coerenza tra le Fasi}

La Phase C \`e il punto di arrivo di un percorso iniziato con l'Ethnographic Assessment (Phase A) e proseguito con il Feasibility Study (Phase B). Ogni decisione progettuale \`e tracciabile alle evidenze raccolte:

\begin{itemize}
    \item Dai test utente $\rightarrow$ alle raccomandazioni R1--R13 $\rightarrow$ all'Information Architecture $\rightarrow$ al Blueprint $\rightarrow$ ai Wireframe.
    \item Le Issue (D1--D5, R1--R4) identificate nell'analisi euristica hanno trovato risposta in una o pi\`u raccomandazioni.
    \item Le Personas non sono rimaste un esercizio accademico: hanno guidato scelte concrete di profondit\`a IA, dimensione del testo, linguaggio e modalit\`a di navigazione.
\end{itemize}

\subsection{Decisioni Critiche}

Alcune decisioni hanno richiesto un bilanciamento tra esigenze contrastanti:

\begin{itemize}
    \item \textbf{Guida ai canali vs. libert\`a di scelta}: Il wizard ``Dove Comprare'' guida l'utente verso il canale ottimale, ma alcuni utenti esperti (come Chris) potrebbero preferire la scelta libera. La soluzione \`e mostrare il canale consigliato in evidenza, ma elencare sempre tutte le alternative in una sezione ``Altri modi per acquistare''.
    \item \textbf{Hub informativo vs. contenuti aziendali}: Mantenere \texttt{tper.it} come sito puramente informativo significa che le news aziendali e i comunicati stampa devono trovare spazio senza oscurare le funzioni principali. La scelta \`e stata di relegarli in una sezione separata ``Azienda'' raggiungibile dal footer, mentre tutta l'area principale \`e dedicata all'utente.
    \item \textbf{Completezza informativa vs. semplicit\`a}: La descrizione di 5+ canali con tutte le loro sfumature rischia di riprodurre il problema originale (sovraccarico informativo). Il compromesso \`e la tabella comparativa filtrabile: l'utente vede tutto ma pu\`o restringere il campo per tipo di operazione o metodo di pagamento.
    \item \textbf{Redirect espliciti vs. esperienza fluida}: Mostrare un avviso prima di reindirizzare a solweb aggiunge un passaggio, ma riduce il disorientamento del cambio dominio. Il test con Chris ha dimostrato che il costo in termini di click aggiuntivo \`e trascurabile rispetto al beneficio di sapere dove si sta andando.
\end{itemize}

\subsection{Contributo delle Personas al Design}

Il contributo pi\`u significativo delle personas \`e stato quello di rendere le decisioni progettuali ``personali''. Quando si doveva scegliere tra due opzioni, la domanda ``Cosa farebbe Maria?'' o ``Cosa capirebbe Fatima?'' ha fornito una bussola inequivocabile. Questo ha garantito che il design rimanesse centrato sull'utente anche nelle scelte tecniche pi\`u astratte.

In particolare, il tema della \textbf{confusione sui canali} --- emerso trasversalmente in tutte e quattro le personas --- \`e diventato il fulcro del redesign informativo di \texttt{tper.it}. Non si trattava di migliorare l'acquisto (che avviene altrove), ma di fare da ``ponte informativo'' tra l'utente e il canale giusto. Questo ha ridefinito la missione stessa del sito: non pi\`u un contenitore generico di informazioni, ma una \textbf{guida attiva all'ecosistema TPER}.

\subsection{Limiti e Sviluppi Futuri}

\begin{itemize}
    \item I wireframe presentati sono in formato lo-fi. Una prototipazione ad alta fedelt\`a con test utente summativi \`e il naturale passo successivo (Phase D).
    \item La guida ai canali e il wizard ``Dove Comprare'' sono progettati per essere testati con utenti reali (Maria, Laura, Chris, Fatima) per validare l'efficacia nel ridurre la confusione di canale.
    \item Il ponte QR code tra Punto Tper e \texttt{tper.it} richiede un'integrazione tecnica con il sistema gestionale dello sportello fisico. La fattibilit\`a tecnica \`e stata verificata in Phase B, ma resta da implementare.
    \item I cross-link ``Intendevi...?'' richiedono un motore di riconoscimento delle intenzioni basato sulla navigazione dell'utente. Una prima versione pu\`o essere implementata con regole statiche, ma l'evoluzione naturale \`e l'uso di machine learning per suggerimenti predittivi.
\end{itemize}

\vspace{0.5cm}

\begin{center}
{\color{secondaryColor}\small
Documento che descrive l'impatto delle personas e degli use case sul design proposto \\[4pt]
\textbf{Phase C --- Design Proposal: Final Reflections} \\[4pt]
Progetto: UX Design Project: Help People Help People --- TPER S.p.A.}
\end{center}

\end{document}
